% FILE: Diagnostic and assessment tools — test protocols, evaluation instruments, assessment methods
\chapter{Diagnostic Tools and Assessment Scales}
\label{app:diagnostic-tools}

This appendix describes the principal diagnostic instruments and assessment tools used in ME/CFS clinical evaluation and research. For each tool, we specify its purpose, administration procedure, scoring, interpretation, and known limitations in ME/CFS populations.

\section{Symptom Questionnaires}
\label{sec:symptom-questionnaires}

\subsection{DePaul Symptom Questionnaire (DSQ)}

The DePaul Symptom Questionnaire is a validated self-report instrument developed specifically for ME/CFS research and clinical assessment~\cite{Jason2006}.

\begin{description}
\item[Purpose] Systematic assessment of ME/CFS symptom frequency and severity across multiple domains; enables classification according to multiple case definitions (Fukuda, CCC, ICC).
\item[Administration] Self-administered paper or electronic questionnaire; approximately 20--30 minutes to complete.
\item[Content] 54 symptom items, each rated on two scales: frequency (0 = none of the time, 4 = all of the time) and severity (0 = symptom absent, 4 = very severe). Additional items capture demographics, illness onset, and functional status.
\item[Scoring] Composite domain scores are calculated for fatigue, post-exertional malaise, sleep dysfunction, pain, neurocognitive symptoms, autonomic dysfunction, neuroendocrine symptoms, and immune symptoms. A symptom is considered clinically significant when frequency $\geq 2$ and severity $\geq 2$.
\item[Interpretation] The DSQ enables operationalized application of diagnostic criteria. For the ICC, PEM must score $\geq 2$ on both frequency and severity, plus at least one symptom in each required category (neurological, immune, energy metabolism).
\item[Strengths] Validated psychometric properties; enables systematic case identification; widely used in ME/CFS research, facilitating cross-study comparison.
\item[Limitations] Self-report introduces subjective variability; cognitively impaired patients may require assistance; does not capture symptom timing or PEM delay characteristics.
\end{description}

\subsection{Bell Disability Scale}

The Bell Disability Scale provides a single-number rating of functional capacity~\cite{Roberts2016}.

\begin{description}
\item[Purpose] Rapid assessment of overall functional impairment; tracking disease severity over time.
\item[Administration] Clinician-rated or self-rated; takes less than 5 minutes.
\item[Scoring] A single score from 0 to 100 in increments of 10:
\begin{itemize}
    \item \textbf{100}: Normal activity, no symptoms
    \item \textbf{80}: Normal activity with effort, mild symptoms
    \item \textbf{60}: Can work part-time with rest periods; moderate symptoms
    \item \textbf{50}: Can perform light duties with rest; unable to work full-time
    \item \textbf{40}: Mostly homebound; severe symptoms with activity
    \item \textbf{30}: Mostly bedridden; severe symptoms at rest
    \item \textbf{20}: Bedridden most of the day; unable to perform self-care independently
    \item \textbf{10}: Bedridden and dependent; severe symptoms continuously
    \item \textbf{0}: Bedridden, unable to sit up; critically ill
\end{itemize}
\item[Interpretation] Scores correspond approximately to NICE 2021 severity grades~\cite{NICE2021mecfs}: Mild ($\geq 60$), Moderate (40--50), Severe (20--30), Very Severe ($\leq 10$).
\item[Strengths] Quick to administer; provides a standardized severity metric; sensitive to change over time.
\item[Limitations] Coarse granularity (10-point increments); does not capture symptom domains; subjective calibration varies between raters.
\end{description}

\section{Functional Assessments}
\label{sec:functional-assessments}

\subsection{SF-36 (Short Form Health Survey)}

The SF-36 is a generic health-related quality of life instrument widely used across chronic diseases, providing a benchmark for comparing ME/CFS functional impact with other conditions~\cite{kingdon2018functional}.

\begin{description}
\item[Purpose] Assessment of health-related quality of life across eight domains; comparison with population norms and other chronic disease groups.
\item[Domains] Physical Functioning (PF), Role-Physical (RP), Bodily Pain (BP), General Health (GH), Vitality (VT), Social Functioning (SF), Role-Emotional (RE), Mental Health (MH).
\item[Scoring] Each domain is scored 0--100, with higher scores indicating better health. Summary scores: Physical Component Summary (PCS) and Mental Component Summary (MCS).
\item[ME/CFS-specific findings] ME/CFS patients score significantly lower than population norms on all domains, with the most severe impairment in RP, PF, SF, and VT. ME/CFS SF-36 scores are comparable to or lower than scores for congestive heart failure, multiple sclerosis, and type~2 diabetes~\cite{kingdon2018functional,hvidberg2015quality}.
\item[Limitations] Generic instrument not designed for ME/CFS; does not capture PEM, orthostatic intolerance, or cognitive symptoms specifically; ceiling effects may limit sensitivity in severely affected patients who are unable to complete the questionnaire.
\end{description}

\subsection{Chalder Fatigue Scale}

The Chalder Fatigue Scale is an 11-item self-report measure of fatigue severity.

\begin{description}
\item[Purpose] Quantification of physical and mental fatigue severity.
\item[Scoring] 11 items scored on a 4-point Likert scale (bimodal scoring: 0, 0, 1, 1) or a 4-point continuous scale (0--3). Total score ranges from 0 to 11 (bimodal) or 0 to 33 (Likert).
\item[Limitations] Developed as a general fatigue measure, not specific to ME/CFS. Does not distinguish ME/CFS-type fatigue from fatigue due to depression, sleep disorders, or other causes. Used as a primary outcome in the controversial PACE trial~\cite{White2011pace}, which has been criticized for using recovery thresholds that overlapped with entry criteria~\cite{geraghty2019cognitive}. Not recommended as a sole outcome measure in ME/CFS trials.
\end{description}

\section{Objective Tests}
\label{sec:objective-tests}

\subsection{Two-Day Cardiopulmonary Exercise Testing (2-Day CPET)}

Two-day CPET is the gold-standard objective test for demonstrating PEM and is the only test that can objectively quantify the exercise intolerance characteristic of ME/CFS~\cite{keller2024cpet}.

\begin{description}
\item[Purpose] Objective demonstration of post-exertional physiological impairment; documentation of reduced exercise capacity for disability evaluation; research quantification of metabolic dysfunction.
\item[Protocol] Two maximal graded exercise tests on consecutive days, typically on a cycle ergometer:
\begin{enumerate}
    \item Day~1: Standard incremental exercise test to volitional exhaustion. Measures peak VO$_2$, anaerobic threshold (AT/VT1), peak workload, and respiratory exchange ratio.
    \item Day~2: Identical protocol repeated 24 hours later.
\end{enumerate}
\item[Interpretation] In healthy controls and most other chronic diseases, day-2 performance equals or exceeds day-1 performance. In ME/CFS, day-2 performance is significantly reduced:
\begin{itemize}
    \item Peak VO$_2$ decreases by $\geq 8$\% (clinically significant decline)
    \item Anaerobic threshold decreases, indicating earlier metabolic shift to glycolysis
    \item Peak workload and exercise duration decrease
\end{itemize}
\item[Clinical significance] The day-2 decrement is objective, reproducible, and specific to ME/CFS among studied conditions. Provides strong evidence for PEM in disability evaluations.
\item[Limitations] Maximal exercise testing is contraindicated in some severely affected patients; provokes PEM lasting days to weeks; limited availability; requires specialized equipment and trained personnel. Not suitable for serial monitoring due to the PEM it induces.
\end{description}

\subsection{Home-Based Assessment Protocols for Severe ME/CFS}
\label{sec:home-based-assessments}

Standard clinic-based diagnostic procedures are poorly tolerated by severe and very severe ME/CFS patients (Bell score $\leq$30), with $>$60\% unable to complete inpatient protocols due to PEM~\cite{Fricke2026achtsam}. The ACHTSAM study~\cite{Fricke2026achtsam} is the first protocol to systematically evaluate which assessments can be safely administered through home visits, using a staged approach that orders procedures from low to high burden with mandatory rest periods.

\begin{description}
\item[Low-burden assessments] ($\leq$5\,min, supine-tolerant): Resting ECG, pupillography, osteosonography (bone density), bioelectrical impedance body composition, near-infrared spectroscopy (NIRS) for muscle oxygenation, light activity sensor.
\item[Moderate-burden assessments] (5--15\,min): HRV with paced breathing, Schellong orthostatic test, handgrip dynamometry, blood sampling, salivary cortisol.
\item[High-burden assessments] (25--30\,min, mental/physical exertion): Cognitive testing (MoCA, TMT-B, SDMT), EndoPAT endothelial function (peripheral arterial tonometry), neurocognitive/psychological assessment.
\item[Design principles] Assessments proceed in burden order; supine procedures precede seated/standing tests; each procedure includes post-assessment perceived exertion rating; patients may pause or terminate at any time.
\item[Tolerability hypothesis] Non-invasive supine-position assessments are expected to achieve $\geq$80\% completion, while high-burden tests requiring sustained mental exertion or standing are expected to achieve $<$30\%.
\item[Status] Protocol published; results expected mid-2026.
\end{description}

\subsection{Tilt Table Testing}

Tilt table testing evaluates autonomic cardiovascular regulation during orthostatic stress.

\begin{description}
\item[Purpose] Diagnosis of orthostatic intolerance, postural orthostatic tachycardia syndrome (POTS), and neurally mediated hypotension.
\item[Protocol] The patient rests supine for 10--20 minutes, then is passively tilted to 60--70$^\circ$ for up to 45 minutes. Continuous heart rate and blood pressure monitoring throughout.
\item[Diagnostic criteria]
\begin{itemize}
    \item \textbf{POTS}: Sustained heart rate increase $\geq 30$\,bpm (or $\geq 40$\,bpm in ages 12--19) within 10 minutes of tilt, without orthostatic hypotension
    \item \textbf{Neurally mediated hypotension}: Drop in systolic BP $\geq 20$\,mmHg or diastolic BP $\geq 10$\,mmHg, often with paradoxical bradycardia
    \item \textbf{Orthostatic hypotension}: Sustained drop in systolic BP $\geq 20$\,mmHg or diastolic $\geq 10$\,mmHg within 3 minutes of tilt
\end{itemize}
\item[ME/CFS relevance] Orthostatic intolerance is present in 50--90\% of ME/CFS patients. Tilt table testing provides objective documentation and identifies the specific subtype, guiding treatment selection.
\item[Alternative] Active standing test (10-minute stand with serial HR and BP measurement) provides a simpler office-based screen, though less standardized.
\end{description}

\subsection{Autonomic Function Tests}

Additional autonomic assessments complement tilt table testing.

\begin{description}
\item[Heart rate variability (HRV)] Analysis of beat-to-beat heart rate intervals from ECG or heart rate monitor recordings. Time-domain (SDNN, RMSSD) and frequency-domain (LF, HF, LF/HF ratio) metrics quantify sympathovagal balance. ME/CFS patients typically show reduced HRV, reduced HF power (parasympathetic), and elevated LF/HF ratio (sympathetic predominance)~\cite{Newton2007}.
\item[Valsalva maneuver] Forced expiration against a closed glottis (40\,mmHg for 15 seconds). The heart rate and blood pressure response has four phases testing baroreflex integrity. Abnormal responses indicate baroreflex failure.
\item[Deep breathing test] Heart rate response to deep breathing at 6 breaths/minute. The expiratory--inspiratory HR difference (E:I ratio) quantifies parasympathetic function. Reduced in ME/CFS.
\end{description}

\section{Laboratory Tests}
\label{sec:laboratory-tests}

ME/CFS is currently a clinical diagnosis without a pathognomonic laboratory test. Laboratory testing serves two purposes: (1)~excluding alternative diagnoses and (2)~identifying treatable comorbidities.

\subsection{Standard Exclusion Panel}

The following tests are recommended to exclude conditions that can mimic ME/CFS~\cite{NICE2021mecfs,Carruthers2011ICC}:

\begin{itemize}
    \item \textbf{Complete blood count}: Excludes anemia, leukemia, infection
    \item \textbf{Comprehensive metabolic panel}: Excludes renal/hepatic disease, electrolyte disorders
    \item \textbf{Thyroid function} (TSH, free T4): Excludes hypothyroidism/hyperthyroidism
    \item \textbf{Erythrocyte sedimentation rate (ESR) and C-reactive protein (CRP)}: Screens for active inflammatory/autoimmune conditions
    \item \textbf{Ferritin}: Excludes iron deficiency (common comorbidity)
    \item \textbf{Vitamin B12 and folate}: Excludes deficiency-related fatigue
    \item \textbf{Cortisol} (morning or ACTH stimulation test): Excludes Addison's disease (primary adrenal insufficiency)
    \item \textbf{Celiac serology}: Excludes celiac disease
    \item \textbf{Urinalysis}: Screens for renal disease, diabetes
\end{itemize}

\begin{warning}[Morning Cortisol Alone Does Not Exclude Post-Viral Secondary Adrenal Insufficiency]
Standard morning cortisol screening is designed to detect primary adrenal insufficiency (Addison's disease), where cortisol is markedly low and ACTH is elevated. Post-viral secondary adrenal insufficiency---documented in $\sim$40\% of SARS-1 survivors~\cite{Leow2005sars} and in COVID-19 pituitary case series~\cite{Carosi2024hypopituitarism}---presents differently: cortisol falls in the low-normal range (typically not flagged as abnormal by laboratory reference ranges), and ACTH is low or inappropriately normal rather than elevated. This pattern is invisible to morning cortisol screening alone. In any post-COVID ME/CFS patient with fatigue disproportionate to other findings, unexplained hypotension, or a clinical picture consistent with cortisol insufficiency, dynamic testing should be considered: ACTH stimulation test (most accessible) or insulin tolerance test (ITT, gold standard for central AI). See Chapter~\ref{ch:endocrine}, Section~\ref{sec:postviral-pituitary} for mechanistic context.
\end{warning}

\subsection{Extended Panel for Clinical Characterization}

Additional tests characterize disease mechanisms but are not required for diagnosis:

\begin{itemize}
    \item \textbf{Immunoglobulin levels} (IgG, IgA, IgM): Identifies humoral immune deficiency
    \item \textbf{Lymphocyte subsets} (CD4, CD8, NK cells): Characterizes immune cell populations
    \item \textbf{NK cell function assay}: Quantifies cytotoxic activity (research setting)
    \item \textbf{Autoantibody panel} ($\beta_2$-adrenergic, muscarinic M3/M4 receptor antibodies): Identifies autoimmune subtype (specialized laboratories)
    \item \textbf{Viral serology} (EBV VCA IgG/IgM, EBNA, HHV-6 IgG): Documents prior/active herpesvirus infection
    \item \textbf{Tryptase}: Screens for mast cell activation
    \item \textbf{Vitamin D}: Common deficiency in housebound patients
    \item \textbf{Sleep study (polysomnography)}: Excludes primary sleep disorders (obstructive sleep apnea, narcolepsy)
\end{itemize}

\subsection{Experimental Biomarkers}

No validated diagnostic biomarker for ME/CFS exists, but several candidates are under investigation:

\begin{itemize}
    \item \textbf{Metabolomics panels}: Reduced amino acids, altered lipid profiles, elevated lactate-to-pyruvate ratios~\cite{Naviaux2016metabolomics,Germain2020metabolic}
    \item \textbf{Cytokine panels}: Disease-duration-dependent profiles~\cite{Hornig2015,Montoya2017}
    \item \textbf{Epigenetic markers}: DNA methylation signatures at specific CpG sites~\cite{deVega2021dna_methylation}
    \item \textbf{Two-day CPET}: While not a blood test, the day-2 decrement serves as a functional biomarker (Section~\ref{sec:objective-tests})
\end{itemize}

\begin{warning}[No Single Diagnostic Test Exists]
ME/CFS remains a clinical diagnosis based on symptom patterns (post-exertional malaise as the cardinal feature) after exclusion of alternative explanations. No laboratory test, imaging study, or biomarker has been validated for routine diagnostic use. Patients and clinicians should be cautious of commercial tests marketed as ME/CFS diagnostics that have not undergone rigorous validation.
\end{warning}
